% 用户内容入口：正文与图表样例
% 建议将后续章节按相同模式继续拆分为 chapter2.tex、chapter3.tex 等文件。

\def\JOUBodyFirstRunningHeader{1 绪论}
\def\JOUBodyFirstTitleCN{1 绪论}
\def\JOUBodyFirstTitleEN{1 Introduction}
\def\JOUBodyFirstSectionTitle{1.1 概述（Introduction）}
\def\JOUBodyFirstDotsTop{……\\……\\……}
\def\JOUBodyFirstSubsectionTitle{1.1.1 研究目标}
\long\def\JOUBodyFirstParagraph{%
\hspace*{2em}描述旋流－静态微泡浮选柱的旋流场结构，分析旋流场特征及其影响；借助\\
流体力学软件对柱体的内部流场进行模拟并分析其流场速度分布规律，研究循环\\
矿浆量及给矿量等因素对流场的影响；通过对旋流场内的颗粒受力分析，建立基\\
于旋流的颗粒动力学方程；系统揭示旋流分选作用，并进行相关动力学分析…%
}
\def\JOUBodyFirstDotsBottom{……\\……}
\def\JOUBodyFirstPageNumber{1}

\def\JOUBodySecondRunningHeader{硕士学位论文}
\def\JOUBodySecondSectionTitle{1.1.2 研究方法}
\def\JOUBodySecondLead{流场模拟及分选机理研究。}
\long\def\JOUBodySecondIntroParagraph{%
为保证旋流分选机理分析的可解释性，本文采用理论分析、实验对比与图表展示相结合的方式组织样例内容。本页保留表格、插图与中英文标题的典型版式，用于后续补充实际论文章节内容。%
}
\def\JOUBodySecondTableTitleCN{筛分粒度组成}
\def\JOUBodySecondTableTitleEN{Particle size distribution results}
\long\def\JOUBodySecondTableRows{%
粒级，mm & 产率，％ & 灰分，％ & 累计产率，％ & 累计灰分，％ \\
\midrule
＞0.5 & 3.80 & 7.38 & 3.80 & 7.38 \\
0.5～0.25 & 4.55 & 4.56 & 8.35 & 5.84 \\
0.25～0.125 & 3.32 & 5.47 & 11.67 & 5.74 \\
0.125～0.074 & 4.74 & 3.63 & 16.41 & 5.13 \\
0.074～0.045 & 10.72 & 3.11 & 27.13 & 4.33 \\
＜0.045 & 72.87 & 4.64 & 100.00 & 4.56 \\
合计 & 100.00 & 4.56 & － & － \\
}
\def\JOUBodySecondSeriesOne{(600,0.040) (800,0.055) (1000,0.070) (1100,0.089) (1250,0.102) (1350,0.108)}
\def\JOUBodySecondSeriesTwo{(600,0.042) (800,0.055) (1000,0.073) (1100,0.092) (1250,0.108) (1350,0.115)}
\def\JOUBodySecondSeriesThree{(600,0.043) (800,0.055) (1000,0.075) (1100,0.094) (1250,0.110) (1350,0.134)}
\def\JOUBodySecondLegendOne{背压为0.37m水柱}
\def\JOUBodySecondLegendTwo{背压为0.90m水柱}
\def\JOUBodySecondLegendThree{背压为1.50m水柱}
\def\JOUBodySecondFigureTitleCN{循环矿浆压力与柱体背压的关系}
\def\JOUBodySecondFigureTitleEN{Relationship between the pressure of circulating pulp and the back pressure}
\def\JOUBodySecondPageNumber{2}
